% Chapter qt, part C: Hecke family, computing the descent, assessment, open problems. Sources: Q lines 634 to 937 and L lines 196 to 346 (misc/reviews/nick-qrh_chapter_for_vadim, nick-levy_mixture_chapter, read as text only), referee reports misc/research/directions/opus/vol3/QRH2.md (M5, M7, M8, M10, M11, M12, M13 and minors) and LEVY.md (M6 to M10), numbers misc/vol3/verify/QRH/results_descent.md, assessment misc/research/rh2/R2.md, R3.md, R4.md (sections 4 and 6); written 20 Nov 2026.
\section{The Hecke family over $\mathbb Q(\sqrt{-3})$}\label{qt:sec:hecke}

Write $F=\mathbb Q(\sqrt{-3})$. For a finite order Hecke character $\eta$ of $F$ with conductor $\mathfrak f_\eta$, complete the $L$ function with the complex Gamma factor,
\[
\Lambda_F(s,\eta)=\big(3\,N\mathfrak f_\eta\big)^{s/2}(2\pi)^{-s}\Gamma(s)\,L_F(s,\eta).
\]
For the trivial character the displayed completion has poles at $0$ and $1$, and one multiplies by $s(s-1)$ before forming the reciprocal clock: for $7/8<\alpha<1$ the reciprocal of the unstripped completion has a zero at the positive real $u=(1-\alpha)^2$, which a probability Laplace transform cannot have. By class field theory $L_F(s,\eta)$ is a degree two $L$ function over $\mathbb Q$. The floor $7/8$ of \cite{OAI26a} covers every finite order Hecke $L$ function over $F$, so for every such $\eta$ the Thorin ladder of $\Lambda_F(s,\eta)$ has the Poisson form of Corollary~\ref{qt:cor:floor} for all $\alpha>7/8$, with every Cauchy width at least $\alpha-\frac78$, uniformly across the family. When $\eta$ is real, $\Lambda_F(s,\eta)$ is self dual and real on the real axis. The pairing lemma (Lemma~\ref{qt:lem:pair}) uses only self duality and the functional equation, so it holds for the family unchanged, and so does the first inequality of Theorem~\ref{qt:thm:C}. The counting step of that lower bound and Theorem~\ref{qt:thm:D} use an explicit zero count for $\zeta$ from \cite{HSW22}, and for $\Lambda_F(s,\eta)$, which has degree two and conductor $3N\mathfrak f_\eta$, they need an explicit count of its own, which we do not have. For complex $\eta$ the normalised single reciprocal can be complex on $u>0$, so one uses a real entire completion from the outset, $\Lambda_F(s,\eta)\Lambda_F(s,\bar\eta)$ with any poles removed, in which the zeros of $\eta$ pair with those of $\bar\eta$. For a primitive unitary character the functional equation combined with conjugation gives the reflection $\rho\mapsto1-\bar\rho$, and reality gives the conjugation symmetry of the real clock. A conductor independent lower bound on the widths supplies no conductor independent positive margin for the compensated minimum.

For the Dedekind zeta function $\zeta_F(s)=\zeta(s)L(s,\chi_{-3})$ the phases add. The completed product is $\xi(s)\Lambda(s,\chi_{-3})$, where $\xi$ carries the factor $s(s-1)/2$, and
\[
U_\alpha^F=U_\alpha^{\zeta}+U_\alpha^{\chi_{-3}},
\]
which is the Thorin form of the factorisation on a common admissible level, generalised Gamma convolutions being closed under convolution. The first zero of $L(s,\chi_{-3})$ is at ordinate $8.0397$, and the density of the product has a Cauchy component centred there on top of those of $\xi$.
% FIGURE: phase densities at level 0.55 for xi, for Lambda(s,chi_{-3}) and for their product, the completed Dedekind zeta function; first extra peak at 8.04.

Windowed positivity extends as far as verification allows. Platt \cite{Platt16} verifies the generalised Riemann hypothesis for primitive Dirichlet characters of modulus $q\le400{,}000$ to height at least $10^8/q$. % CHECK Platt16: exact height statement, and that the modulus of chi chi_{-3} stays within the range
For the Dedekind zeta function of $F$ the binding factor is $L(s,\chi_{-3})$ with $q=3$, so the zeros of $\zeta_F$ are known to lie on the critical line up to ordinate about $3.3\cdot10^7$. For a symmetric zero set verified through $H$ the elementary uncompensated conclusion of Theorem~\ref{qt:thm:window} holds through $H-\frac12$. The compensated and curvature statements of Part B used a zeta specific baseline, low gaps and a zero count, and need analogues proved for the new function before they are claimed on that window. For a norm pullback $\eta=\chi\circ N_{F/\mathbb Q}$ with $\chi$ primitive, $L_F(s,\eta)=L(s,\chi)L(s,\chi\chi_{-3})$ up to the Euler factors at primes dividing the conductors, and the factors must be taken with their primitive inducing characters: $\chi=\chi_{-3}$ pulls back to the trivial primitive character, and treating $\chi_{-3}^2$ as the principal character modulo $3$ would insert a spurious factor $1-3^{-s}$ into $\zeta_F$. Platt's range applies only when both primitive conductors are in it, with the smaller of the two heights. For ray class characters that do not come from $\mathbb Q$ no family wide verification is known, but there are exceptions. Let $K=\mathbb Q(\sqrt[3]2)$, $F=\mathbb Q(\sqrt{-3})$ and $M=KF$; the splitting field has group $S_3$, and either nontrivial character $\eta$ of $\mathrm{Gal}(M/F)=C_3$ defines a finite order Hecke character of $F$, exchanged with its inverse by conjugation, whereas a norm pullback is invariant. Induction of $\eta$ is the irreducible degree two summand of the permutation representation on the three embeddings of $K$, so $\zeta_K=\zeta\,L_F(s,\eta)$. Tollis verified rigorously the cubic field of discriminant $-108$ to height $92$, and by the factorisation and the holomorphy of $L_F(s,\eta)$ away from the real pole at $1$ its zeta verification gives the same finite window for this non pullback $L_F$. The family wide problem stays open, and where no factorisation through a verified field is available only the levels above $7/8$ are unconditional.

For a complex character the L\'evy weights $2\Re\,\eta(\mathfrak a)\Lambda_F(\mathfrak a)$ are signed already for $\alpha>1$, so the $2\sin^2$ lemma does not give compensated positivity by itself, while the uncompensated density is positive there by the level criterion.
% FIGURE: Thorin density, compensated density and bare Euler product phase density for the product of a cubic Hecke character and its conjugate at level 1.5 (say whether it is a Dirichlet character modulo 7 or the cubic residue character over F).

\section{Computing the descent}\label{qt:sec:compute}

The descent of the floor chapter, from a level where the phase measure is known to a lower one, is a deconvolution by Cauchy kernels, and positivity of the measure is the only structural information it can use (Proposition~\ref{bo:prop:descent}, Theorem~\ref{tf:thm:deconv}). Under RH, $\theta_1'=\pi\,P_{1/2}*\sum_\gamma\delta_\gamma$ on the whole line. With RH verified only through $T_0$ the exact equation is $v_1(y)/\pi=\sum_{|\gamma|\le T_0}P_{1/2}(y-\gamma)+\sum_{|\gamma|>T_0}P_{1-\beta}(y-\gamma)$, and the second term is not determined by finite RH verification; for $|y|\le W<T_0$ the tail is bounded by $\pi^{-1}\sum_{|\gamma|>T_0}(|\gamma|-W)^{-2}$, and the corresponding all on line tail is bounded similarly. Finite positive fits do not certify the zero set, so the computations below are reconstruction tests. Given the whole line equation with all widths $\frac12$, one solves once for a positive measure $x$ at the centre and climbs to any level by $\theta_\alpha'=\pi\,P_{\alpha-1/2}*x$. Here $P_\varepsilon(d)=\varepsilon/(\pi(\varepsilon^2+d^2))$ is the Cauchy kernel. The data are $b=\theta_1'/\pi$ on $|y|\le70$ with grid step $0.05$, and the unknown lives on the wider grid $|t|\le100$, so that zeros outside the data window are modelled; padding the grid lets exterior mass affect the fit but does not absorb it exactly, and a truncation at $100$ or $120$ introduces a separate, much nearer tail that also needs a model or a bound. With $K$ the discretised kernel $P_{1/2}$ we compare the following iterations.

\begin{enumerate}
\item EM, the multiplicative iteration $x\leftarrow x\cdot K^\top(b/Kx)/K^\top\mathbf 1$ of Richardson \cite{Richardson72} and Lucy \cite{Lucy74}, put in likelihood form for positive linear inverse problems by Vardi and Lee \cite{VardiLee93}. For positive $K$ and $b$, a positive start and nonzero column sums it keeps every iterate positive and the Poisson likelihood is nondecreasing; the observed decay of its Kullback Leibler gap is like $1/k$ (a proved $O(1/k)$ bound on the objective suboptimality needs assumptions that are not verified here, a raw Kullback Leibler residual may have a nonzero optimum, and the reconstruction error is measured in a different norm).
\item Nesterov EM: an EM step from the extrapolated point $y_k=x_k+\frac{t_k-1}{t_{k+1}}(x_k-x_{k-1})$ \cite{Nesterov83}, with $t_1=1$ and $t_{k+1}=(1+\sqrt{1+4t_k^2})/2$, kept positive (an extrapolated point with a nonpositive entry is rejected), with a restart from $x_k$ whenever the Kullback Leibler objective rises \cite{OC15}. A rejected step falls back to a plain EM step, so the likelihood is monotone by construction. Section~\ref{cx:sec:accel} uses the geometric form of the same momentum, in $\log x$, where positivity is automatic.
\item Projected gradient and its accelerated form (FISTA with gradient restart) for $\frac12\|Kx-b\|^2$ on $x\ge0$, with no penalty.
\end{enumerate}

The error is the relative $L^1$ distance of $P_{0.1}*x$ from $\theta_{0.6}'/\pi$ on $|y|\le40$. After $50$, $200$, $1000$ and $3000$ iterations it is $0.50$, $0.35$, $0.22$ and $0.155$ for EM, and $0.28$, $0.10$ and about $0.02$ to $0.03$ at both $1000$ and $3000$ iterations for Nesterov EM (reported experiments). % FIGURE: error at level 0.6 against iteration for the three methods, and the recovered centre measure with the zero ordinates marked.
The fitted exponents of the error between iterations $30$ and $1000$ are about $-0.27$ for EM and $-0.77$ for Nesterov EM, and $-0.21$ for projected gradient against $-0.43$ for FISTA.

A spectral heuristic, for the unconstrained linear problem and only suggestive for EM, accounts for the gain. The Cauchy kernel of width $\varepsilon$ has symbol $e^{-\varepsilon|\omega|}$. A gradient method resolves frequency $\omega$ after about $e^{2\varepsilon\omega}$ iterations and an accelerated one after about $e^{\varepsilon\omega}$, so at a given iteration count acceleration doubles the resolved band. The observed exponents of projected gradient against FISTA are in the ratio about two, and those of EM against Nesterov EM in the ratio about three. The Kullback Leibler gap of EM decays like $1/k$, and that of Nesterov EM like about $k^{-1.7}$ between $k=100$ and $1000$ against a reference limit from a run of $15{,}000$ iterations. We have no proof of any rate for Nesterov EM with the Cauchy kernel; Remark~\ref{qt:rem:accelerated} gives a $k^{-2}$ bound for a different accelerated method.

\begin{remark}[An accelerated entropic mirror descent has a $k^{-2}$ bound]\label{qt:rem:accelerated}
Write the objective as $G(x)=\sum_i\{b_i\log\frac{b_i}{(Kx)_i}-b_i+(Kx)_i\}$, $B=\sum b_i$, $c_j=\sum_iK_{ij}$, $Q_{ij}=K_{ij}/c_j$, $d_i=b_i/B$, $s=\sum_jc_jx_j$ and $p_j=c_jx_j/s$. Then $G(x)=BF(p)+s-B-B\log(s/B)$ with $F(p)=\sum_id_i\log\{d_i/(Qp)_i\}$, so the observed mass is profiled out and $G_*=B\min_{\Delta}F$. Because each row of the column normalised Cauchy matrix has a finite ratio $M_i/m_i$ of its largest to its smallest entry, $F$ has a gradient that is Lipschitz in the $\ell^1$ norm on the simplex with constant $L_0=\tfrac14\sum_id_i\{(M_i-m_i)/m_i\}^2$, and the entropy is strongly convex in $\ell^1$ (Pinsker). In this setting Nesterov's accelerated method with the entropy prox function on the simplex, % TODO bib: Nesterov 2005, Smooth minimization of non smooth functions, Math. Program. 103, 127 to 152
that is, the accelerated entropic mirror descent with three sequences, step $1/(L\theta_k)$ and $\theta_k=2/(k+2)$, which may take an EM step at every iteration and keep the better of the two, has true objective gap $G(x_N)-G_*\le4BL_0D(p_*\|p_0)/(N+1)^2$, with $D$ the Kullback Leibler divergence. This is an upper bound and not an attained order, and for the Cauchy kernel $M_i/m_i\le(1+T/\varepsilon)^4$ on a source window of diameter $T$, uniformly in the observation sites and the grid.

The method is not the Nesterov EM of the experiments (one EM step from a clipped extrapolation, with restart), for which no rate is proved; the rate comes entirely from the auxiliary sequence. And the constants are too large to explain the observed rate: on the chapter's grid $T/\varepsilon=400$ and the uniform bound for $L_0$ is of order $10^{20}$, and the bound at $N=1000$ is far above the observed gap of a few $10^{-3}$ (in the units of $G$).
\end{remark}

\subsection{Scale mixtures and moving atoms}\label{qt:sec:mix}

The grid method treats the mixing measure as weights on fixed cells. A faster method makes the positions of the atoms the parameters and uses the Cauchy kernel as a scale mixture of normals, in the form that Section~\ref{cx:sec:mixture} develops. By the hierarchical duality of Polson and Scott \cite{PolsonScott16} the latent scale $\lambda$ of the mixture also gives an envelope: $\log(1+u)=\min_{\lambda>0}\{\lambda(1+u)-\log\lambda-1\}$, attained at $\lambda=1/(1+u)$, so that
\[
-\log P_\varepsilon(d)=\min_{\lambda>0}\Big\{\lambda\,\frac{\varepsilon^2+d^2}{\varepsilon^2}-\log\lambda\Big\}+\text{const},\qquad\hat\lambda(d)=\frac{\varepsilon^2}{\varepsilon^2+d^2}.
\]
For fixed $\lambda$ the bound is quadratic in the position, so the position step is a weighted mean. Model the measure by atoms $\sum_kw_k\delta_{t_k}$ with widths $e_k$, write $m(y)=\sum_kw_kP_{e_k}(y-t_k)+F(y)$ with a fixed far field $F$, and use the Poisson likelihood on the window $|y|\le W$. With $d_{ik}=y_i-t_k$ and grid step $h$, one iteration is
\begin{align*}
n_{ik}&=h\,b(y_i)\,w_kP_{e_k}(d_{ik})/m(y_i),\\
\lambda_{ik}&=e_k^2/(e_k^2+d_{ik}^2),\qquad w_k=\sum_in_{ik}/c_k,\\
t_k^{\rm new}&=\frac{\sum_in_{ik}\lambda_{ik}\,y_i-\frac12 e_k^2\,w_k\big(P_{e_k}(W+t_k^{\rm old})-P_{e_k}(W-t_k^{\rm old})\big)}{\sum_in_{ik}\lambda_{ik}},
\end{align*}
where $c_k$ is the mass of atom $k$ inside the window. The second term in the numerator accounts for the tails outside the window. Without it the iteration converges quickly to a biased answer: the error from level $1$ stalls near $1.5\cdot10^{-4}$ and it grows with the width, but we claim no power law for the bias. The window term is evaluated at the old position $t_k^{\rm old}$, so the update is lagged, the step is not an exact majorise minimise step and monotonicity of the likelihood is not guaranteed.

For a single Cauchy atom of known width observed as population data on the whole line, the position step $M$ has the true position as a fixed point and satisfies $M(t)-t_*=\frac12(t-t_*)$ for every $t$ and every width. This is Proposition~\ref{cx:prop:ratehalf}, where $\frac12$ is identified with the fraction of missing information of the location EM for a Cauchy law. The doubled step $t+2(M(t)-t)$ therefore returns $t_*$ in one iteration. For atoms that overlap, that is when the width $\alpha-\frac12$ is comparable to the spacing of the zeros, the rate degrades, and a lattice model says by how much.

\begin{proposition}[Rate for a lattice]\label{qt:prop:lattice}
For an infinite periodic lattice of unit atoms at spacing $\Delta$ with known width $\varepsilon$ and weights, observed as population data on the whole line, the linearised population position step at the true lattice has the rigid translation of all atoms (a Bloch mode) as a mode, with contraction factor
\[
\kappa=1-\varepsilon^2\int_0^\Delta\frac{m'(y)^2}{m(y)}\,dy,\qquad m(y)=\frac{\sinh(2\pi\varepsilon/\Delta)}{\Delta\,(\cosh(2\pi\varepsilon/\Delta)-\cos(2\pi y/\Delta))},
\]
and, with $A=2\pi\varepsilon/\Delta$, one has exactly $1-\kappa=A^2/(2\sinh^2A)$; for all sufficiently large $A$ it is the slowest mode, and
\[
1-\kappa\sim8\pi^2\Big(\frac\varepsilon\Delta\Big)^2e^{-4\pi\varepsilon/\Delta}.
\]
\end{proposition}

\begin{proof}
The Jacobian of an EM map at the truth is $I-I_c^{-1}I_o$, with complete and observed information, and the same form holds for the position block of the majorise minimise step, since by symmetry the weights decouple from the rigid mode. The complete information is $I_c=\varepsilon^{-2}$ per atom: conditional on the envelope precision the location curvature is $2\lambda/\varepsilon^2$, and $E\lambda=\frac12$ at the truth. The observed information is diagonalised by the Bloch modes $e^{ik\varphi}$, with eigenvalue $\int_0^\Delta\big|\sum_ke^{ik\varphi}P_\varepsilon'(y-k\Delta)\big|^2/m(y)\,dy$, and the rigid observed information per cell is $\int_0^\Delta(m')^2/m$, since at $\varphi=0$ the sum is $-m'$; this gives $\kappa$. For the exact formula put $c=\cosh A$ and $u=2\pi y/\Delta$: the mean of $\sin^2u/(c-\cos u)^3$ over a period is $1/[2(c^2-1)^{3/2}]$, by differentiating the mean $1/\sqrt{c^2-1}$ of $1/(c-\cos u)$, which gives $1-\kappa=A^2/(2\sinh^2A)$ and checks both constants in the asymptotic. For the asymptotics write $a=2\pi\varepsilon/\Delta$ and $u=2\pi y/\Delta$, so that $m=\Delta^{-1}\big(1+2\sum_{n\ge1}e^{-na}\cos nu\big)$. As $a\to\infty$, $m\to\Delta^{-1}$ and $m'\to-2\Delta^{-1}(2\pi/\Delta)e^{-a}\sin u$, so $m'^2/m\to16\pi^2\Delta^{-3}e^{-2a}\sin^2u$, whose integral over a period is $8\pi^2\Delta^{-2}e^{-2a}$, and multiplication by $\varepsilon^2$ gives the stated form. For minimality put $t=|\varphi|/(2\pi)\in[0,\frac12]$, $q=e^{-A}$ and $b_n=(n+t)q^{|n+t|}$. Fourier expansion of the Bloch derivative and $1/m\propto\cosh A-\cos u$ show that the observed information, up to a positive factor independent of $t$, is $N(t)=(1+q^2)\sum_nb_n^2-2q\sum_nb_nb_{n+1}$. Completing the squares outwards on $n\ge0$ and $n\le-1$ leaves $b_0^2+b_{-1}^2-2qb_0b_{-1}$, and the geometric sums give, with $C_q=q^2/(1-q^2)$, $N(t)=(t^2+C_q)e^{-2At}+((1-t)^2+C_q)e^{-2A(1-t)}+2q^2t(1-t)$. Here $N(0)=2C_q$ and $N'(0)=0$. On $0\le t\le1/(8A)$ the first term has second derivative at least $e^{-1/4}$, since its quadratic multiplier is $2-8At+4A^2(t^2+C_q)$, the second term has nonnegative second derivative there for large $A$, and the last contributes $-4q^2$; hence $N(t)>N(0)$ for $0<t\le1/(8A)$. On the remaining interval $N(t)\ge t^2e^{-2At}\ge e^{-A}/(64A^2)>2C_q$ for sufficiently large $A$. This proves minimal observed information, hence maximal contraction factor, exactly at the rigid mode for all sufficiently large $A$; no claim is made for all widths.
\end{proof}

We have not proved that the rigid mode is the slowest for $\varepsilon/\Delta\le0.3$, the regime in which the iteration is used. The values of $\kappa$ and of the inverse spectral gap $1/(1-\kappa)$ from the closed form are as follows.

\begin{center}
\begin{tabular}{rrr}
$\varepsilon/\Delta$ & $\kappa$ & inverse gap $1/(1-\kappa)$\\
0.05 & 0.51613 & 2.1\\
0.20 & 0.69709 & 3.3\\
0.30 & 0.82835 & 5.8\\
0.50 & 0.96300 & 27.0\\
0.80 & 0.99782 & 459.6\\
1.00 & 0.99972 & 3631.7
\end{tabular}
\end{center}

The asymptotic formula agrees with these values to three digits from $\varepsilon/\Delta=0.8$ on and is not yet valid at $0.3$ to $0.5$. At height $T$ the mean spacing of the zeros is $2\pi/\log(T/2\pi)$, so $\varepsilon/\Delta=(\alpha-\frac12)\log(T/2\pi)/(2\pi)$, and the gap is $1-\kappa\sim8\pi^2(\varepsilon/\Delta)^2(T/2\pi)^{-2(\alpha-1/2)}$.
% FIGURE: the gap 1 - kappa against epsilon over Delta, with the asymptotic form.

\subsection{Experiments}\label{qt:sec:exp}

In the reported experiments at low height the data are $b$ on $|y|\le70$ with step $0.05$ at levels $\alpha=1$, $1.5$ and $2$, and the iteration is started from the grid solution (300 iterations of Nesterov EM, the clusters of mass inside $|t|<66$ taken as atoms), whose error is of order $10^{-3}$. The zeros with $66<|\gamma|<120$ are supplied at their true positions with free weights, and beyond $120$ the far field is the integral of $P_\varepsilon$ against the smooth zero density $\frac1{2\pi}\log\frac t{2\pi}$. The error is the largest distance from a true zero below $45$ to its atom, over the eight zeros below $45$. From level one the mixture step reaches an error of about $10^{-5}$ in about ten iterations, and the doubled step gets there in fewer, with first step ratios of about $0.55$, close to $\frac12$. From levels $1.5$ and $2$, where the data are sums over prime powers and the iteration is a map from primes to zeros, the best errors are of order $10^{-5}$ to $10^{-4}$ and $10^{-3}$, and their exact values depend on how the clusters of the grid stage are extracted, so we quote only the order. Nesterov momentum on top of the mixture step does not help: from iteration $20$ on it is worse at every level.
% FIGURE: error against iteration at the three levels for the mixture step, the doubled step and the momentum variant; grid Nesterov EM from a flat start in grey.

If the widths are left free the same iteration estimates $\alpha-\beta$ for each zero, since a zero at $\beta$ has an atom of width $\alpha-\beta$ at level $\alpha$, provided $\alpha>\beta$ (a zero at the level contributes an atom and one to its right a signed term) and the model carries one atom for each zero at an ordinate; distinct coincident components of the same width are identifiable only by their total weight. Here the widths of the exterior atoms (the supplied zeros beyond the core) are held fixed and only the core atoms are free. When the exterior widths are also free the estimates are biased, with widths $0.37$ to $0.45$ at level one. Starting from the true positions, with $200$ iterations, the widths of the zeros below $45$ are within $0.01$ of $\alpha-\frac12$ at the three levels, with weights within $0.01$ of one. It certifies nothing, since the zeros below $45$ are known to be on the line, the comparison uses supplied exterior ordinates and true starting positions, and no control with a planted pair $\frac12\pm b+i\gamma$, or with the Davenport and Heilbronn function \cite{Davenport,Spira94}, was run, so that identifiability of a width different from $\alpha-\frac12$ is not shown. The far field assumes the width $\alpha-\frac12$. At level two a width of $1.49$ could not by itself tell $\beta=\frac12$ from $\beta=0.51$.

At greater height the zeros crowd and the same level is a harder problem. Take the window $|y-10^4|\le30$, which contains $70$ zeros, with mean spacing $0.853$. The atoms of the fit cover $|t|<27$ about the centre, $63$ zeros. Then $\varepsilon/\Delta=0.29$ at $\alpha=0.75$ and $0.59$ at $\alpha=1$, and the lattice factors are $\kappa=0.82$ and $0.98$. From $\alpha=0.75$ the grid stage recovers $62$ of the $63$ atoms, and the mixture step then reaches errors of order $10^{-5}$ to $10^{-4}$, slowly, in hundreds of iterations. From level one the grid stage merges close pairs, a split of the clusters of mass near $k$ into $k$ atoms assumes simple zeros and is a hypothesis of the method, and the draft results for level one at height $10^4$  are not reproduced by an independent implementation and are withdrawn. What is supported is an empirical operating range, not a resolution theorem: descend from a level with $\varepsilon/\Delta$ at most about $0.3$, that is $\alpha-\frac12\le1.9/\log(T/2\pi)$, which at $T=10^4$ allows $\alpha-\frac12$ up to $0.26$. Level one satisfies it only up to heights near $270$, and levels above one, where the data are prime sums, only below that.

\section{The hitting time strip}\label{qt:sec:strip}

The reciprocal xi law at level $\alpha\ge\tfrac12$ is the law of the random variable $X_\alpha\ge0$ with $\mathbb E\,e^{-sX_\alpha}=L_\alpha(s)=\xi(\alpha)/\xi(\alpha+\sqrt{2s})$, which exists for every real $\alpha\ge\tfrac12$: at the centre by Theorem~\ref{thm:ch12:idsd}, and at shifted levels by Theorem~\ref{bp:thm:alpha}, where $X_\alpha$ is a Brownian first passage time to a random level (in the book's normalisation $G_\alpha(u)=L_\alpha(u/2)$ the clock is $X_\alpha/2$). Let $B$ be a standard Brownian motion from $0$. This section asks for a boundary $g$ such that the exit time $\tau_g=\inf\{t>0:|B_t|\ge g(t)\}$ of the strip $|x|<g(t)$ has the law of $X_\alpha$, the inverse first passage problem of Anulova \cite{Anulova80}. By Brownian scaling, in the book's chart $T=X_\alpha/2$ the boundary is $g_\alpha(2T)/\sqrt2$. The one sided problem, a barrier $b$ for $B_t\ge b(t)$, is that of Ekstrom and Janson \cite{EkstromJanson16}. We write $\ell=\ell(t)>1$ for the large root of $\ell-\log\ell=\log\frac1{4\pi t}$ ($0<t<1/(4\pi e)$), so that $\ell=\log\frac1t+\log\log\frac1t-\log4\pi+o(1)$.

\begin{theorem}[The strip exists exactly for $\alpha\ge\tfrac12$]\label{qt:thm:stripexist}
For every real $\alpha\ge\tfrac12$ there is a lower semicontinuous symmetric boundary $g_\alpha$ with $\tau_{g_\alpha}\sim X_\alpha$; it is unique in that class, finite, positive and continuous on $(0,\infty)$, and $g_{\alpha_n}\to g_\alpha$ uniformly on compact subsets of $(0,\infty)$ when $\alpha_n\to\alpha$. For $\alpha<\tfrac12$ no law with transform $L_\alpha$ exists. The law is a generalised gamma convolution exactly when no zero has $\Re\rho>\alpha$ (Theorem~\ref{bo:thm:level}), for $\alpha\ge\tfrac78$ by \cite{OAI26a}, and unconditionally for $\alpha\ge1$.
\end{theorem}

Existence of the strip therefore carries no information about RH at any level $\alpha\ge\tfrac12$. The sketch is as follows. At the centre the law of $X_{1/2}$ is infinitely divisible with a positive L\'evy density (Theorem~\ref{thm:ch12:idsd}); for $\alpha>\tfrac12$, $L_\alpha(s)=L_{1/2}(\tfrac{a^2}2+s+a\sqrt{2s})/L_{1/2}(\tfrac{a^2}2)$ with $a=\alpha-\tfrac12$ uses only values of $L_{1/2}$ on the positive axis, so no zero location enters (Theorem~\ref{bp:thm:alpha}). Anulova's theorem then gives the lower semicontinuous symmetric boundary, survival is positive at every finite time, and uniqueness is the theorem of Klump and Kolb, continuity in the level is that of Klump and Savov, together with a modulus of continuity from a uniform lower bound for the densities. Below the centre $\xi(\tfrac12+x)=\int\Phi(u)e^{xu}du$ with a positive kernel, so $\log\xi$ is strictly convex and $b_\alpha=(\log\xi)'(\alpha)<0$ for $\alpha<\tfrac12$, whence $L_\alpha(s)=1-b_\alpha\sqrt{2s}+O(s)>1$ for small $s>0$, impossible for a law on $[0,\infty)$. % TODO bib: Klump and Kolb, symmetric inverse first passage problem, Theorems 2.1 and 2.2; Klump and Savov, Proposition 3.2; theorem numbers to be checked against the papers


\section{Is there a path to RH?}\label{qt:sec:path}

We see no path through the Thorin ladder alone, and the computations above cannot change that: they are a verification method on a window, limited by the amplification $(T/2\pi)^{\alpha-1/2}$ at the lattice frequency and by the precision of the data, not a proof. The reasons below are limitations of the arguments we have tried and not impossibility theorems, and they sharpen what the book states elsewhere.

First, a rung below the floor is a zero free statement only in the following precise senses. A GGC proof for the full reciprocal at one level $\alpha$ proves $\Theta\le\alpha$. Nonnegative phase density, or nonnegative compensated excess, on every level above $\alpha_0$ gives $\Theta\le\alpha_0$ (Corollary~\ref{pol:cor:equiv} for the compensated version, Theorem~\ref{qt:thm:ggc} and Remark~\ref{rem:ch2:oneline} for the other distinctions), and an argument that reads a rung off the prime sums converges only above the abscissa $\Theta$ of the Chebyshev error (Proposition~\ref{ct:prop:qrh}). The floor of \cite{OAI26a} is $\frac78$, and a positivity proof on every line to the right of $\alpha$ yields the half plane $\re s>\alpha$, which for $\alpha>\frac78$ is already known. Positivity at a regular height or on a verified finite window gives a local conclusion only, as Theorem~\ref{qt:thm:window} shows, and the finite reconstructions in this chapter establish no new zero free half plane.

Second, the one sided bound constrains the outcome of a proof and cannot preclude a new one sided method. The one sided Riesz mean bound of Proposition~\ref{qt:prop:riesz} is equivalent to the two sided bound at the same exponent, to the zero free half plane at that exponent and to positivity of the ladder above it, so it is not a relaxation that a positivity argument could exploit, and positivity of $\Lambda$ plays no role in the equivalence. In its RH form the inequality is attained in the limit (the slack of the positive definiteness statement has $\liminf S=0$), which is what an honest target must do. The Davenport and Heilbronn function has a first zero off the line, $0.808517+85.699348i$ \cite{Spira94}, and by Rodgers and Tao \cite{RodgersTao} the inequality $\Lambda_{\rm DN}\ge0$ implies a nonreal zero of the standard $H_t$ for every $t<0$, which after the conversion of the variable and symmetry is a zero off the line. These are controls for individual proof mechanisms: a mechanism that applies uniformly along the heat flow reaches at best an exponent $\Theta_t>\frac12$, and a mechanism that uses the Euler product applies to neither. We claim no impossibility theorem from these controls.

Third, level one is where the arithmetic positivity of the chain ends, in the following sense. For $\alpha>1$ one has the transform identity $G^\infty_\alpha(u)=G_\alpha(u)\,\zeta(\alpha+\sqrt u)/\zeta(\alpha)$ with $G^\infty_\alpha=A(\alpha)/A(\alpha+\sqrt u)$, the archimedean clock. The first factor on the right is a GGC. The arithmetic factor is the transform of an infinitely divisible first passage law (in the standard Brownian convention the first passage level is $(\log M_\alpha)/\sqrt2$ for $\Pr(M_\alpha=n)=n^{-\alpha}/\zeta(\alpha)$); its nonconstant cosine series density has mean zero and changes sign, and it is positive at ordinate zero in the sign convention of this transform, while a bare Euler factor in the reciprocal completion has the opposite sign. The archimedean clock and the arithmetic clock lie on opposite sides of this identity, and the archimedean clock is not a GGC for $1<\alpha<\alpha_*\approx2.660659$ by a numerical scan, not by a proved interval. The positivity of the product is therefore a property of the whole of it. Below level one the renormalised level measure $-e^{-\alpha x}d(\psi(e^x)-e^x)$ is signed and the factorisation is the budget identity of the Chebyshev chapter, and we know of no arithmetic positive object there.

Finite Euler products and the rates are treated in Remark~\ref{qt:rem:eulerlevel}.

\begin{proposition}[Level one without zeros]\label{qt:prop:levelone}
\begin{enumerate}
\item For every $\sigma>0$ and real $t$, $|\xi(1-\sigma+it)|<|\xi(1+\sigma+it)|$. Equivalently $S(z)=\xi(z)/\xi(1+z)$ is holomorphic on $\Re z>0$ with $|S|<1$ there.
\item Consequently, in the notation of Theorem~\ref{h:thm:crit} at $\alpha=1$, the chain of level one is regular: $m_0=\xi'(1)/\xi(1)>0$ and every $a_j(1)>0$, and both families of leading Hankel determinants are positive.
\end{enumerate}
\end{proposition}

This is already in the book, through the Hadamard product: a pair of zeros contributes a factor larger than $1$ to the quotient as soon as the offset from the line is at most the level shift, which is Proposition~\ref{fm:prop:bruijn} and the proof of Theorem~\ref{cm:thm:tilt}, the positivity of the Hadamard sum used in Proposition~\ref{pol:prop:round}, and Theorem~\ref{h:thm:crit}(i) at $\alpha=1$. What is new is a proof that uses no zero of $\xi$ at all: only the Euler product on $\Re s>1$, the functional equation and a crude polynomial bound for $\zeta$.

\begin{proof}
(1) By $\xi(s)=\xi(1-s)$ and conjugation, $|\xi(1-\sigma+it)|=|\xi(\sigma+it)|$, so the claim is $|S(z)|<1$ for $z=\sigma+it$, and $\xi(1+z)\ne0$ there by the Euler product. For $0<\varepsilon<\tfrac12$ put $S_\varepsilon(z)=\xi(z-\varepsilon)/\xi(1+\varepsilon+z)$, holomorphic on $\Re z>-\varepsilon$, with $|S_\varepsilon(it)|=1$ because $\xi(it-\varepsilon)=\overline{\xi(1+\varepsilon+it)}$. Partial summation gives $\zeta(v)=\tfrac1{v-1}+\tfrac12-v(v+1)\int_1^\infty P(u)u^{-v-2}du$ with $P(u)=(\{u\}^2-\{u\})/2$ bounded, so $|\zeta(v)|\ll(1+|v|)^2$ on $\Re v\ge-\tfrac12$. With $|1/\zeta(1+\varepsilon+z)|\le\zeta(1+\varepsilon)$ on $\Re z\ge0$ and Stirling's formula for the gamma ratio (of size $|z|^{-1/2-\varepsilon}$), $|S_\varepsilon(z)|\le C_\varepsilon(1+|z|)^2$ on the closed half plane. The maximum principle on half discs for $S_\varepsilon(z)(1+\delta z)^{-3}$, then $\delta\to0$ and $\varepsilon\to0$, gives $|S|\le1$. For real $x\to\infty$, $S(x)\sim\sqrt{2\pi/x}\to0$, so $S$ is not a unimodular constant and the inequality is strict.

(2) By $F(-z)/F(z)=S(z)$ for $F(z)=\xi(1+z)/\xi(1)=h(z^2)+zg(z^2)$, the Cayley transform $H_0=(1-S)/(1+S)=(F(z)-F(-z))/(F(z)+F(-z))=z\varphi(z^2)$, $\varphi=g/h$, is odd and holomorphic near $0$, real on the positive axis, has $\Re H_0>0$ on $\Re z>0$ by (1), and $H_0(x)\to1$. Let $H$ be any function with these properties and $H(x)\to\ell\in(0,\infty)$, and $c=H'(0)$. The Schwarz and Pick inequality at a positive real point $x$, with the squared moduli subtracted from one, gives $H(x)\Re H(z)/|H(z)+H(x)|^2\ge x\Re z/|z+x|^2$, and $x\downarrow0$ gives $c\,\Re(1/H)\ge\Re(1/z)>0$, so $c>0$. Then $H_*=c/H-1/z$ is odd and holomorphic near $0$, real on the axis, has $\Re H_*\ge0$, tends to $c/\ell>0$ on the axis, so $\Re H_*>0$ by the minimum principle, and the same reduction applies again. By induction every $c_j=H_j'(0)>0$, where $H_{j+1}=c_j/H_j-1/z$, and with $H_j=z\varphi_j(z^2)$ this reads $\varphi_j(w)=c_j/(1+w\varphi_{j+1}(w))$, which is the Stieltjes elimination: $c_0=m_0$ and $c_j=a_j(1)$. The determinants are $D_n^{(0)}=c_0^n\prod_{j=1}^{n-1}(c_{2j-1}c_{2j})^{n-j}$ and $D_n^{(1)}=D_n^{(0)}\prod_{j=1}^nc_{2j-1}$, both positive.
\end{proof}

The extension to a level $\alpha>\tfrac12$ is not proved and is the ladder criterion: a zero $\beta+i\gamma$ with $\beta>\alpha$ makes the right side of $|\xi(\alpha-\sigma+it)|<|\xi(\alpha+\sigma+it)|$ vanish at $\sigma=\beta-\alpha$, $t=\gamma$, so the strict inequality at every such centre would imply RH.

Our assessment is that there are three live targets. The compensated density $Z_\alpha+P_\alpha+G^\Gamma_\alpha>0$ on the strip, stated entirely in the language of the book, whose first order form on the critical line, the curvature inequality, is proved on the window. A one line form of the level one rung, positivity of $\re\xi'/\xi(1+it)$ without the Hadamard product, by the explicit constant $\sum1/\rho=0.0231$ at the origin and the bounds for $\zeta'/\zeta(1+it)$ at large height, the template of which at level $1-\eta$ names exactly the strip estimate needed next. And the control test on any claimed proof of a one sided bound (see the problems below). We do not expect the descent algorithm to be more than a check on these.

\section{Open problems}\label{qt:sec:open}

Related problems posed elsewhere are cited by their number in Section~\ref{fin:sec:open}: the smallest detectable displacement of a zero from $\frac78$ by the EM descent (item 44), the existence and attainment of a least level $\alpha_P$ at which the truncated Euler product is a generalised Gamma convolution and its rate of approach to $1$ (item 43; that $\inf\{\alpha:T^{(P)}_\alpha\text{ is a GGC}\}\to1$ itself follows from Remark~\ref{qt:rem:eulerlevel}, since for every $\epsilon>0$ all large $P$ work at $1+\epsilon$ and no $P$ works below $1$), the acceleration of the EM descent with momentum at heights $10^4$ to $10^6$ (item 45), positivity of the rung $\alpha=1$ for $\zeta$ through the rates from the Euler product (item 17), and the items 53 to 62 on the quasi Riemann hypothesis.

\begin{enumerate}
\item\label{qt:op:rate} Prove a rate for Nesterov EM with the Cauchy kernel. An accelerated entropic mirror descent has the rate $k^{-2}$ for the profiled objective (Remark~\ref{qt:rem:accelerated}), so what is open is the rate of the EM form of the momentum.
\item\label{qt:op:lattice} Prove that the rigid translation is the slowest mode of the position step of Proposition~\ref{qt:prop:lattice} for every $\varepsilon/\Delta$, and not only for all sufficiently large $A=2\pi\varepsilon/\Delta$, and give the rate for irregular spacing, where the closest pair of zeros and not the mean spacing may set the rate.
\item\label{qt:op:stop} Give a stopping rule against the drift of the mixture step and a treatment of the window edge and far field that does not move the minimum of the error, together with a principled split and merge step for the number of atoms at height.
\item\label{qt:op:control} Run the free width estimate with a planted pair $\frac12\pm b+i\gamma$ and with the Davenport and Heilbronn function as controls, and determine the smallest $b$ it detects at a given level and height.
\item\label{qt:op:hecke} Give explicit zero counts for $\Lambda_F(s,\eta)$ of degree two, so that the counting step of the lower bound and the hidden zero theorem hold for the Hecke family, and verify the generalised Riemann hypothesis to modest height for the family of finite order Hecke characters of $\mathbb Q(\sqrt{-3})$ as a whole.
\item\label{qt:op:shadow} Make the shadow of a zero uniform in height. With $S(\alpha,y)$ as in Section~\ref{qt:sec:shadow}, $S\ge0$ at every regular point of the line $\re s=\alpha$ is monotonicity of the compensated density in the level on that line. A weaker form is proved (Proposition~\ref{qt:prop:shadow}, with a two sided quantitative form for a given zero): if $\frac12<\alpha\le1$ and $S(\alpha,\cdot)\ge0$ at every regular point of the line, then $\xi$ has no zero with $|\sigma-\alpha|\le(\alpha-\frac12)/\big(8M(T)\log(8M(T)/(\alpha-\frac12))\big)$ at height $T$, where $M(T)=10+\log(1+|T|)$. What stays open is the sharper width: prove that the same hypothesis forces no zeros with $|\beta-\alpha|<c\,(\beta-\frac12)\min\{1,(\log\gamma)^{-1/2}\}$, handling clusters through the fact that each zero contributes signed mass $\pi\operatorname{sgn}(\alpha-\beta)$ to the $S$ term (a reflected pair has mass $2\pi$ only when $\alpha-\frac12$ exceeds its offset from the line). The numerics give $c\approx0.37$ at low height.
\item\label{qt:op:epstein} Apply the control test to a claimed proof of a one sided Riesz mean bound by theta lifts: find the step that produces an inequality for $\sum_{n\le x}\Lambda(n)(x/n-1)$ and apply it to the Epstein zeta function of $x^2+5y^2$, which has a functional equation and zeros off the line and in $\re s>1$ but no Euler product. If the step goes through there the claim is wrong, and if it fails it identifies an additional hypothesis that must be checked.
\end{enumerate}
